% Part: incompleteness
% Chapter: theories-computability
% Section: introduction

\documentclass[../../../include/open-logic-section]{subfiles}

\begin{document}

\olfileid[pt]{inc}{tcp}{int}

\olsection{Introdução}

\begin{editorial}
Esta seção deve ser reescrita.
\end{editorial}

Temos o seguinte:
\begin{enumerate}
\item uma definição do que significa uma função ser representável em
  $\Th{Q}$ (\olref[req][int]{defn:representable-fn});
\item uma definição do que significa uma relação ser representável em
  $\Th{Q}$ (\olref[req][rel]{defn:representing-relations});
\item um teorema que afirma que as funções representáveis de $\Th{Q}$
  são exatamente as computáveis
  (\olref[req][int]{thm:representable-iff-comp})
\item um teorema que afirma que as relações representáveis de $\Th{Q}$
  são exatamente as computáveis
  \olref[req][rel]{thm:representing-rels})
\end{enumerate}

Uma {\em teoria} é um conjunto de sentenças fechado dedutivamente,
isto é, com a propriedade de que, sempre que $T$ prova $!A$, então
$!A$ pertence a $T$. Talvez seja melhor pensar em uma teoria como uma
coleção de sentenças juntamente com tudo o que elas implicam. De agora
em diante, usaremos $\Th{Q}$ para designar a {\em teoria} constituída
pelo conjunto de sentenças deriváveis dos oito axiomas em
\olref[req][int]{sec}. Lembre-se de que podemos codificar fórmulas de
$\Th{Q}$ como números; se $!A$ é uma dessas fórmulas, seja $\Gn{!A}$ o
número que codifica $!A$. Mediante essa codificação, podemos agora
perguntar se diversos conjuntos de fórmulas são ou não computáveis.

\end{document}
